\begin{statement}{C.2.NC.2}{Proved}{The normalized double curves and the coherent Euler formula}
Let $S$ be a connected compact normal-crossings divisor in a smooth compact complex threefold $Y$ with $H^2(Y,\Q)=0$. Write $\nu:\coprod_iT_i\to S$ for normalization, $D_i\subset T_i$ for its conductor, $C_\alpha$ for the connected components of the normalization of the double locus, and $\tau$ for the number of triple points. Let $R_\alpha$ be the reduced divisor on $C_\alpha$ of preimages of triple points, counted separately for the three double-curve branches at each triple point. Then
\[
 e(C_\alpha)=\deg R_\alpha,\qquad
 \sum_\alpha e(C_\alpha)=3\tau,
 \qquad
 \sum_i\chi(T_i,\OO_{T_i})=\frac{\tau}{2}.
 \tag{C.NC.4}
\]
In particular every normalized double curve has genus at most one, and $\tau$ is even. The remaining topological identities are
\[
\begin{aligned}
 \sum_i e(D_i)&=2\sum_\alpha e(C_\alpha)-3\tau,\\
 e(S)&=\sum_i e(T_i)-\sum_\alpha e(C_\alpha)+\tau.
\end{aligned}
 \tag{C.NC.5}
\]
\end{statement}
\begin{proof}
Adjunction and the local residue comparison give a canonical holomorphic line-bundle isomorphism
\[
 \nu^*\omega_S|_{T_i}\simeq\omega_{T_i}(D_i),
 \qquad \omega_S=(\omega_Y\otimes\OO_Y(S))|_S.
 \tag{C.NC.6}
\]
For example, the residue of $dx\wedge dy\wedge dz/(xy)$ on the branch $x=0$ is $\pm dy\wedge dz/y$, and at $xyz=0$ it is $\pm dy\wedge dz/(yz)$. These expressions agree under changes of local equations by the usual determinant and unit factors of the displayed line bundles, proving the comparison. Since both line bundles on the right in the expression for $\omega_S$ come from $Y$, its first Chern class and that of $\omega_{T_i}(D_i)$ vanish rationally.

Normalize every irreducible curve in every $D_i$. The resulting disjoint union $\widehat D$ maps to $\coprod_\alpha C_\alpha$ by an étale map of total degree two over each $C_\alpha$. Away from a triple point this follows from the two branches of $xy=0$. At a triple point, each of the three normalized double branches has exactly two preimages, and their maps in the remaining coordinate are local isomorphisms. This proves étaleness at the missing points as well. The covering over a given $C_\alpha$ may be connected; no splitting is assumed.

Adjunction for a nodal divisor on a smooth surface, pulled back to $\widehat D$, identifies the restriction of $\omega_{T_i}(D_i)$ with the canonical line of that normalized conductor curve twisted by all its preimages of nodes of $D_i$. Such nodes are exactly the preimages of triple points of $S$. This is the elementary local dualizing-form identity $dw/w$ on each branch of a node. The resulting degree is zero by (C.NC.6). Since the covering of $C_\alpha$ is étale and its marked divisor is the pullback of $R_\alpha$, it follows that
$\deg\omega_{C_\alpha}+\deg R_\alpha=0$.
Thus $e(C_\alpha)=\deg R_\alpha$. There are exactly three marks over every triple point, even when two or three belong to one $C_\alpha$, proving their summed formula.

For the coherent calculation, the alternating branch-resolution sequence is
\[
 0\longrightarrow\OO_S\longrightarrow\nu_*\OO_{\coprod T_i}
 \longrightarrow\bigoplus_\alpha(j_\alpha)_*\varepsilon_\alpha
 \longrightarrow\bigoplus_{p\in S^{[3]}}\C_p
 \longrightarrow0.
 \tag{C.NC.7}
\]
Here $j_\alpha:C_\alpha\to S$ is the finite normalization map of that double curve, and $\varepsilon_\alpha$ is the sign line of its two-sheeted conductor covering. Locally order the branches and send $(f_i)$ to $(f_j-f_i)_{i<j}$; the next map is the alternating sum of the three values at a triple point. The sequence is exact at $xy=0$ because a pair of functions glues precisely when its restrictions agree, and at $xyz=0$ because pairwise differences satisfy exactly the single alternating relation. Equivalently, compare coefficients of convergent power series on the three coordinate planes; equality on their coordinate axes is exactly the condition for these series to come from one series modulo $xyz$. Changing a local branch ordering multiplies the appropriate difference by $-1$, producing $\varepsilon_\alpha$. At a triple point this sign convention only chooses a basis of the one-dimensional last term. Hence the sequence is globally well-defined with the indicated local systems, including self-intersections.

The sign line satisfies $\varepsilon_\alpha^{\otimes2}\simeq\OO_{C_\alpha}$, so it has degree zero and $\chi(C_\alpha,\varepsilon_\alpha)=1-g(C_\alpha)=e(C_\alpha)/2$. The divisor sequence on $Y$ and Hirzebruch--Riemann--Roch give $\chi(S,\OO_S)=0$, since $c_1(\OO_Y(S))=0$ rationally. Taking Euler characteristics in (C.NC.7) now gives
\[
 0=\sum_i\chi(T_i,\OO_{T_i})-\frac12\sum_\alpha e(C_\alpha)+\tau,
\]
which proves the last equality in (C.NC.4).

Finally $\widehat D$ has Euler number $2\sum_\alpha e(C_\alpha)$. Each triple point produces three nodes in $\coprod D_i$, and identifying the two preimages of each node subtracts one. This proves the first identity in (C.NC.5). The double locus in $S$ has normalization $\coprod C_\alpha$ and three preimages over each triple point, so its Euler number is $\sum_\alpha e(C_\alpha)-2\tau$. Normalization is an isomorphism off the conductor. Additivity of Euler characteristic therefore gives the second identity in (C.NC.5).
\end{proof}


% C.2.NC.3 from part_C2_NC.tex

